% Dot batch 71; source SHA-256: bb65562e3cadee226b53026043ecee57e394d148b0b4811f8c4d911435087cec
% Archived source: dot_updates_20261124/batches/batch071/EGA第五卷_定理16.9_光滑二次型的判据.tex
% Original body lines: 19--97; integration map: integration/late_integration.json
\begingroup
\def\newscope{V-Q}
\def\OO{\mathcal O}
\def\Qsp{\mathcal Q}
\def\Sym{\operatorname{Sym}}
\def\Isom{\operatorname{Isom}}
\def\Hom{\operatorname{Hom}}
\egas{16.9}{定理}
设$S$为预概形，$E$为$S$上有限正秩的局部自由模，$Q\in\Gamma(S,\Sym^2 E)$。以下两条等价；若$E$的秩恒为$n$，它们还与(iii)、(iii bis)等价。\cmt{L18-071-01}{原文未显式排除零秩。按15.1的光滑型定义和16.4的$n>0$标准型约定，补写正秩；参见命题16.2的脚注。原文把恒秩限定放在(iii)之内。}
\begin{enumerate}
\renewcommand{\labelenumi}{(\roman{enumi})}
\item $Q$光滑。
\item 校正判别式$\widetilde d(Q)\in\Gamma(S,(\det E)^{\otimes2})$可逆。
\item 存在满射$S'\to S$，使得$Q$经基变换得到的$Q_{S'}$与$n$变量的标准二次型同构。
\item[(iii bis)] 同(iii)，但要求$S'\to S$忠实平坦且为有限表示（fppf）。
\end{enumerate}

\egas{16.10}{推论}
设$k$为域，$E$为$k$上正维有限维向量空间，$Q\in\Sym^2 E$。以下条件等价：
\begin{enumerate}
\renewcommand{\labelenumi}{(\roman{enumi})}
\item $Q$光滑；
\item $\widetilde d(Q)\ne0$；
\item 存在$k$的一个扩域$k'$，使得$Q\otimes_k k'$与标准型同构；
\item[(iii bis)] 同(iii)，且$k'/k$有限。
\end{enumerate}
这直接由16.9推出。注意，当$k$为代数闭域时，(iii)与(iii bis)可换成更醒目的条件：$Q$已经与标准型同构。

\noindent\textbf{定理16.9的证明。}\quad 可以设$E$的秩恒为$n$。由于标准型的校正判别式是单位，\cmt{L18-071-02}{两底本写等于$1$。按16.1和16.5的极化行列式约定，实际为$(-1)^{\lfloor n/2\rfloor}$，仍为单位。}立刻有
\[
  \text{(iii bis)}\Longrightarrow\text{(iii)}\Longrightarrow\text{(ii)} .
\]
现在证明$\text{(ii)}\Rightarrow\text{(iii bis)}$，然后证明$\text{(i)}\Leftrightarrow\text{(ii)}$。

设$E_0=\OO_S^n$，$Q_0$为$E_0$上的标准型。考虑函子
\[
 \Isom((E_0,Q_0),(E,Q)):(\mathrm{Sch}/S)^{\mathrm{op}}\longrightarrow\mathrm{Sets}.
\]
它在每个$S$上的$S'$所取的值，是与$Q_{0,S'}$和$Q_{S'}$相容的同构$(E_0)_{S'}\to E_{S'}$的集合。不必对$E$或$Q$另加条件，这一函子由一个在$S$上仿射且有限表示的预概形$P$表出。它是
\[
 W=\mathbf W\bigl(\Hom_{\OO_S}(E_0,E)\bigr)
\]
的子预概形，也是$W$的开子概形$\Isom(E_0,E)$中的闭子预概形。\cmt{L18-071-03}{两底本在此处把同构空间的第二个模写作$E_1$，据上下文改作$E$。} 要证明$\text{(ii)}\Rightarrow\text{(iii bis)}$，只需证明$\widetilde d(Q)$可逆时$P$在$S$上忠实平坦；事实上可取$S'=P$。

简记$\Qsp(E)=\mathbf W(\Sym^2 E)$，同样定义$\Qsp(E_0)$。通过结构输运，得到
\[
 u\longmapsto\Qsp(u):
 \Isom(E_0,E)\longrightarrow\Isom(\Qsp(E_0),\Qsp(E)).
\]
用$Q_0$对应的截面$q_0$，得到态射
\[
 (*)\qquad \Isom(E_0,E)\longrightarrow\Qsp(E),\qquad
 u\longmapsto\Qsp(u)(q_0).
\]
另一方面，$Q$对应于$\Qsp(E)$在$S$上的截面$q$；按定义，$P$正是$q(S)$在$(*)$下的逆像。这也随即证明前述可表性：$P$在$S$上仿射且有限表示。

在$\Qsp(E)$中引入开子概形$\Qsp(E)^*=\Qsp(E)_{\widetilde d}$，对应于校正判别式可逆的二次型。它表出将$S'$映至$\Sym^2 E_{S'}$中校正判别式可逆的截面集合的函子。结构输运给出
\[
 \Isom(E_0,E)\longrightarrow\Isom(\Qsp(E_0)^*,\Qsp(E)^*).
\]
由于标准型$Q_0$的校正判别式是单位，$q_0$是$\Qsp(E_0)^*$的截面。\cmt{L18-071-04}{法文确为$q_0$、$\Qsp(E_0)^*$；英译混作$q$与$\Qsp(E)^*$，后面又把$\Qsp(E)^*$误排为$O(E)^*$。原稿数值$1$改作单位，理由同上。} 因而$(*)$通过下列态射分解：
\[
 (**)\qquad \Isom(E_0,E)\longrightarrow\Qsp(E)^* .
\]
两端都在$S$上有限表示，所以此态射也有限表示。现在只需证明下述命题。

\egas{16.11}{命题}
沿用以上记号，态射$(**)$忠实平坦。

因此，若$Q$满足(ii)，即$q$实际上是$\Qsp(E)^*$的截面，则沿$q$对$(**)$作基变换所得的$P$仍在$S$上忠实平坦。这就证明了16.9的$\text{(ii)}\Rightarrow\text{(iii bis)}$。

\noindent\textbf{命题16.11的证明。}\quad 两端都在$S$上光滑，所以只需逐纤维证明平坦性，这把问题归结到$S$是代数闭域$k$的谱的情形。这时可以设$E_0=E$，于是$(**)$成为
\[
 (***)\qquad \mathrm{GL}(E)\longrightarrow\Qsp(E)^*,\qquad u\longmapsto u(q_0).
\]
此态射来自群概形$\mathrm{GL}(E)$在$\Qsp(E)^*$上的自然作用，其中$q_0$对应于标准型$Q_0$。由泛平坦性定理，$\Qsp(E)^*$在$k$上光滑，因而约化；存在稠密开集$U\subseteq\Qsp(E)^*$，其上$(***)$平坦。对于任意$u\in\mathrm{GL}(E)(k)$，即$E$的任一自同构，$u(U)$也满足同一平坦性条件。要证明$(***)$平坦，只需证明所有$u(U)$覆盖$\Qsp(E)^*$。为此，只要证明$\mathrm{GL}(E)(k)=\operatorname{Aut}(E)$在$\Qsp(E)^*(k)$上作用传递，即一切校正判别式非零的二次型都在标准型的轨道上。这也给出$(***)$的满射性，故归结为下列引理。\cmt{L18-071-05}{这段末尾原稿只称“判别式”；由$\Qsp(E)^*$的定义，此处应为校正判别式，尤其是在特征2的奇秩情形。}

\egas{16.11.1}{引理}
设$Q$是代数闭域上的$n$变量二次型，且$\widetilde d(Q)\ne0$，则$Q$与标准型同构，亦即适当选取基就能写成标准型。

若$E$为偶秩或$k$的特征不为$2$，假设意味着$Q$非退化；结论可见于 Bourbaki。反之，若$k$的特征为$2$且$E$为奇秩，则$Q$退化。取关联双线性型的核中的一条直线$E_1\subseteq E^\vee$及其补空间$E_2$。于是$Q$写成秩一型$Q_1$与偶秩型$Q_2$的直和$Q_1\oplus Q_2$。\cmt{L18-071-06}{法文及英译在选直线处写在$E$中；按16.1，$B$定义在$E^\vee$上，核与补空间均应取在$E^\vee$中。} 由16.8，$Q_1$、$Q_2$的校正判别式都非零。由前一情形，$Q_2$与标准型同构；又$Q_1$非零，故与一变量标准型同构。于是$Q$与标准型同构，引理得证。

还需证明16.9中(i)与(ii)等价。可以设$S$为代数闭域的谱，从而归结到：
 
\egas{16.11.2}{引理}
$Q$光滑，当且仅当$Q$与标准型同构。\cmt{L18-071-07}{法文手写题句是双向箭头；英译题句只印正向，随后的证明实含两个方向。原稿后述极化箭头有颠倒之处，按16.1统一作$E^\vee\to E$。}

由16.2，可以设$k$的特征为$2$且$E$为奇秩。若$Q$与标准型同构，由16.3可立即验证其光滑：$E^\vee\to E$的核为一维，$Q$在核上不恒为零，因此在核中只有平凡零点。

反过来，仍由16.3，若$Q$光滑，那么$Q$限制到极化根基$N$后仅有平凡零点。因此$N$为一维，且$Q|_N\ne0$。取$N$的一个补空间，便得到$Q$与标准型同构。这证明了16.11.2，完成16.9的证明。
\par\endgroup
